\documentclass[journal,onecolumn]{IEEEtran}

% ── Encabezado ──────────────────────────────────────────────
\usepackage{fancyhdr}
\pagestyle{fancy}
\fancyhf{}
\renewcommand{\headrulewidth}{0.4pt}
\fancyhead[L]{\small Metodología de Ingeniería de Requisitos}
\fancyfoot[C]{\thepage}

% ============================================================
% IDIOMA Y CODIFICACIÓN  (mover ANTES de csquotes)
% ============================================================
\usepackage[utf8]{inputenc}
\usepackage[T1]{fontenc}
\usepackage[spanish,es-tabla]{babel}

% ============================================================
% CITAS Y TEXTO
% ============================================================
\usepackage{csquotes}

% ============================================================
% ESPACIADO Y LISTAS
% ============================================================
\usepackage{parskip}
\setlength{\parskip}{4pt}
\usepackage{enumitem}

% ============================================================
% IDIOMA Y CODIFICACIÓN
% ============================================================
\usepackage[spanish,es-tabla]{babel}
\usepackage[utf8]{inputenc}
\usepackage[T1]{fontenc}

% ============================================================
% PAPEL Y GEOMETRÍA
% ============================================================
\usepackage{geometry}
\geometry{a4paper, top=2.5cm, bottom=2.5cm, left=2.5cm, right=2.5cm}

% ============================================================
% IMÁGENES
% ============================================================
\usepackage{graphicx}
\graphicspath{{Figuras/}}

% ============================================================
% TABLAS
% ============================================================
\usepackage{booktabs}
\usepackage{tabularx}
\usepackage{longtable}

% ============================================================
% MATEMÁTICAS
% ============================================================
\usepackage{amsmath}
\usepackage{amssymb}

% ============================================================
% CÓDIGO FUENTE
% ============================================================
\usepackage{listings}

% ============================================================
% FORMATO ACADÉMICO
% ============================================================
\usepackage{microtype}
\usepackage{float}

% ============================================================
% CAPTIONS (TÍTULOS DE FIGURAS/TABLAS)
% ============================================================
\usepackage{caption}
\captionsetup{font=small, labelfont=bf}

% ============================================================
% REFERENCIAS / HIPERVÍNCULOS
% ============================================================
\usepackage[hidelinks]{hyperref}

% ============================================================
% NUMERACIÓN Y NOMBRES EN ESPAÑOL
% ============================================================
\renewcommand{\thetable}{\arabic{table}}
\def\tablename{Tabla}

\renewcommand{\thesection}{\arabic{section}}
\renewcommand{\thesubsection}{\thesection.\arabic{subsection}}
\renewcommand{\thesubsubsection}{\thesubsection.\arabic{subsubsection}}

% ============================================================
% ESTILO DE TÍTULOS DE SECCIÓN
% ============================================================
\usepackage{titlesec}
\titleformat{\section}{\normalfont\Large\bfseries}{\thesection}{1em}{}
\titleformat{\subsection}{\normalfont\large\bfseries}{\thesubsection}{1em}{}
\titleformat{\subsubsection}{\normalfont\normalsize\bfseries}{\thesubsubsection}{1em}{}

% ============================================================
% CADA SECCIÓN PRINCIPAL (\section) INICIA EN PÁGINA NUEVA
% ============================================================
\usepackage{etoolbox}
\pretocmd{\section}{\clearpage}{}{}

% ============================================================
% BIBLIOGRAFÍA (BIBLATEX + BIBER)
% ============================================================
\usepackage[
backend=biber,
style=ieee,
sorting=none
]{biblatex}
\addbibresource{referencias_metodologia.bib}

% ============================================================
% DOCUMENTO
% ============================================================
\begin{document}

% ============================================================
% SECCIÓN 2: METODOLOGÍA DE INGENIERÍA DE REQUISITOS
% Responsable: GUTIÉRREZ ORTEGA, Génesis Adriana
% ============================================================

\section{Metodología de Ingeniería de Requisitos}
\label{sec:metodologia}

El equipo de la PE5 no se conformó desde la primera práctica. Durante PE1--PE4, cada
integrante trabajó el mismo sistema del PFC --- MundiPets --- en un subgrupo distinto del
curso, y solo a partir de la PE5 el docente reorganizó a los cinco integrantes en un único
equipo. Esta circunstancia no resta validez al proceso: los cuatro subgrupos aplicaron, de
manera independiente, el mismo marco
metodológico sobre el mismo dominio de problema, lo que permite ahora contrastar y
consolidar cuatro ejecuciones reales de un mismo proceso de Ingeniería de Requisitos en
lugar de una sola. La presente sección describe ese proceso tal como se ejecutó,
práctica por práctica, y no un proceso ideal tomado de un manual: cada actividad que se
describe a continuación está respaldada por un artefacto verificable en los repositorios
de PE1--PE4 del equipo.

\subsection{PE1: fundamentos y contexto del sistema}

La primera práctica no consistió en elicitar requisitos de campo, sino en construir el
marco conceptual necesario para hacerlo con criterio en la práctica siguiente. La actividad
central fue doble: por un lado, la caracterización de MundiPets como sistema  su tipo
según la clasificación de sistemas reactivos, transformacionales e interactivos que
propone Pohl \cite{pohl2025}, y la identificación de sus stakeholders mediante una matriz
de poder e interés fundamentada en el PMBOK 7.\textsuperscript{a} edición \cite{pmbok7};
por otro, la aplicación de las ocho perspectivas de contexto de Pohl (dominio, uso,
organización, tiempo, sujeto, desarrollo, intención y tecnología) junto con una declaración
de visión en el formato propuesto por Moore \cite{pohl2025}. 

Se optó por esta técnica y no
por una elicitación directa porque, en ausencia de un marco de análisis previo, una primera
sesión con usuarios corre el riesgo de producir una lista de deseos sin estructura ni
cobertura demostrable de las dimensiones que un sistema sociotécnico como MundiPets
 una plataforma de adopción, cruza responsable y bienestar animal efectivamente
involucra. Los cuatro subgrupos coincidieron en clasificar a MundiPets como un sistema de
información colaborativo con un núcleo interactivo, componentes reactivos (notificaciones)
y un componente transformacional (el cálculo de compatibilidad), y en identificar como
stakeholders de alto interés o alto poder a los dueños de mascotas, los administradores de
la plataforma, las organizaciones de adopción, los veterinarios y los patrocinadores.

El artefacto que produjo la PE1 fue, en consecuencia, un mapa de restricciones y
necesidades preliminares no todavía requisitos formales, sino los temas y las preguntas
que era necesario llevar a la siguiente práctica: qué información sanitaria es relevante
para permitir una interacción entre mascotas, qué datos deben protegerse por ser sensibles,
y qué papel debe jugar la validación veterinaria en el proceso de emparejamiento. Este mapa
es exactamente lo que la PE2 convirtió en una guía de entrevista y en un cuestionario
concretos: sin la perspectiva del dominio y la perspectiva de intención definidas en la
PE1, la PE2 no habría tenido un criterio explícito para decidir qué preguntar a los
propietarios de mascotas ni por qué.

\subsection{PE2: elicitación, análisis y priorización}

Con el marco de la PE1 como entrada, la PE2 ejecutó la elicitación de requisitos
propiamente dicha. La técnica principal fue la \textbf{entrevista estructurada},
complementada con un \textbf{cuestionario digital} y, en algunos subgrupos, con una sesión
de \textbf{brainstorming} grupal para ampliar la cobertura de ideas antes de la
consolidación. La selección de estas dos técnicas como núcleo del proceso y no una
sola respondió a una limitación reconocida por el propio equipo: la entrevista aporta
profundidad cualitativa y permite explorar, en un intercambio directo con el entrevistado,
matices que un formulario cerrado no captura (por ejemplo, qué información considera
privada un propietario y por qué), pero alcanza a pocos participantes y depende de la
disponibilidad real de los stakeholders; el cuestionario, en cambio, permite recoger
información de un número mayor de usuarios potenciales y cuantificar preferencias
mediante escalas Likert, pero no permite repreguntar ni aclarar una respuesta ambigua. 

Usar solo una de las dos habría dejado un vacío: limitarse a la entrevista habría producido una
muestra pequeña y difícil de generalizar; limitarse al cuestionario habría impedido
descubrir necesidades implícitas que los participantes no saben nombrar hasta que se les
pregunta directamente. El IREB CPRE FL sustenta esta selección al describir la entrevista y
el cuestionario como técnicas complementarias cuya combinación reduce el riesgo de sesgo de
una sola fuente \cite{ireb2026}.

La aplicación concreta a MundiPets es trazable a un acta real. En la entrevista
estructurada aplicada a una propietaria de mascota, ante la pregunta de qué información
considera privada o sensible sobre su mascota, la entrevistada señaló la ubicación exacta,
los datos personales del dueño y los detalles médicos que no deberían mostrarse
públicamente sin autorización. Esta respuesta, tomada de forma aislada, no es un requisito;
se convirtió en uno solo después de pasar por el análisis de la PE2, que la clasificó como
una necesidad de privacidad y protección de datos y la tradujo en un requisito no funcional
de confidencialidad de la información sensible del perfil. La misma entrevista, ante la
pregunta sobre qué información es necesaria antes de permitir que una mascota interactúe
con otra, arrojó como respuesta la necesidad de conocer vacunas al día, desparasitación,
enfermedades recientes y antecedentes de comportamiento agresivo: esa respuesta es el
origen directo de los requisitos de validación sanitaria y de compatibilidad que
posteriormente se refinaron en la PE3. La cadena evidencia--necesidad--requisito queda así
documentada, no supuesta.

Una vez recogidos, los requisitos crudos se sometieron a un proceso de análisis,
clasificación (funcional, no funcional, conforme al modelo de calidad de producto de
ISO/IEC 25010 \cite{iso25010}) y priorización. Para priorizar se aplicaron dos técnicas
complementarias y no una sola, cada una con un propósito distinto. \textbf{MoSCoW}
clasificó los requisitos en las categorías Debe tener, Debería tener, Podría tener y No
tendrá esta vez, y resolvió la pregunta de qué construir primero dado un alcance acotado de
tiempo; esta técnica era necesaria porque el equipo identificó más funcionalidades de las
que un MVP inicial podía cubrir, y sin un criterio explícito de priorización la selección
del alcance habría quedado a discreción de cada integrante. 

El \textbf{modelo de Kano}, por
su parte, no compite con MoSCoW sino que responde una pregunta distinta: no cuánto urge un
requisito, sino qué tipo de satisfacción produce en el usuario si se cumple o si falta
 básico, de desempeño o atractivo. Esta distinción resultó relevante para MundiPets
porque algunas funcionalidades, como la validación veterinaria de la información sanitaria,
resultaron básicas (su ausencia genera insatisfacción aunque su presencia no sea percibida
como un plus), mientras que otras, como el cálculo de compatibilidad asistido por
inteligencia artificial, resultaron atractivas (su ausencia no se nota, pero su presencia
sí genera valor diferencial). Aplicar solo MoSCoW habría tratado ambos tipos de requisito
de forma indistinguible; aplicar solo Kano no habría resuelto el problema operativo de
qué construir primero. El artefacto que produjo la PE2 fue la lista depurada y priorizada
de requisitos funcionales, no funcionales y de restricción, con su categoría MoSCoW y su
clasificación Kano, que sirvió de entrada directa a la especificación formal de la PE3.

\subsection{PE3: especificación y modelado}

La PE3 tomó los requisitos priorizados de la PE2 y los llevó a una forma verificable y
modelada. La primera actividad fue la corrección de la redacción de los requisitos
funcionales y no funcionales mediante la \textbf{sintaxis EARS} (Easy Approach to
Requirements Syntax) \cite{mavin2009}, que obliga a expresar cada requisito bajo un patrón
fijo (universal, condicional, opcional o de estado) en lugar de en lenguaje natural
ambiguo. Se eligió EARS porque los requisitos crudos de la PE2, redactados en el lenguaje
de la entrevista y el cuestionario, contenían formulaciones no verificables del tipo
``el sistema debe ser fácil de usar'' que EARS obliga a reescribir con un disparador, una
condición y una respuesta observable del sistema. Esta corrección es la que, más adelante,
hace posible calcular la métrica de verificabilidad (M3) en la sección de métricas de
calidad de este informe: un requisito redactado en EARS es, por construcción, un requisito
verificable.

Sobre los requisitos ya corregidos se construyeron tres modelos complementarios. El
\textbf{modelo de datos} partió de un diagrama entidad--relación que se transformó en un
diagrama de clases UML, siguiendo la práctica de traducir las entidades del dominio
(mascota, propietario, veterinario, solicitud de adopción o cruza) en clases con sus
atributos y relaciones. 

El \textbf{modelo funcional} combinó diagramas de flujo de datos
(DFD) con diagramas de actividades UML para representar cómo se mueve la información a
través de los procesos principales del sistema publicar un perfil, evaluar
compatibilidad, gestionar una solicitud. 

El \textbf{modelo de comportamiento} se resolvió
con una máquina de estados UML (por ejemplo, para el ciclo de vida de una solicitud de
adopción: creada, en revisión, aprobada o rechazada) y con diagramas de secuencia UML que
detallan la interacción temporal entre los actores y el sistema para los casos de uso
principales. La razón de emplear tres modelos y no uno solo es que cada uno responde una
pregunta distinta que ninguno de los otros dos contesta por sí mismo: el modelo de datos
responde qué información existe y cómo se relaciona, el modelo funcional responde cómo
fluye esa información entre procesos, y el modelo de comportamiento responde en qué orden
ocurren los eventos y cómo cambia de estado un objeto del dominio a lo largo del tiempo.

Finalmente, la PE3 cerró con una \textbf{matriz de trazabilidad} que integró los tres
modelos entre sí y con los requisitos de origen, verificando que cada requisito funcional
tuviera al menos un caso de uso, una clase y un elemento del modelo de comportamiento
asociado. Esta matriz preliminar es la base directa de la matriz de trazabilidad final que
se cierra en la sección de trazabilidad y matriz final de este informe: la PE3 no inventó
una estructura nueva para la PE5, sino que extendió la que ya existía.

\newpage

\subsection{PE4: validación mediante inspección formal y gestión de cambios}

Con el ERS modelado en la PE3, la PE4 sometió el documento a una \textbf{inspección formal
	de tipo Fagan} \cite{fagan1976}, reconocida por ISO/IEC/IEEE 29148:2018 entre las técnicas
admisibles del proceso de validación de requisitos \cite{ieee29148}. Se seleccionó la
inspección formal en lugar de una revisión informal porque el método de Fagan exige
preparación individual documentada antes de la reunión, roles diferenciados (Moderador,
Lector, Inspector 1 e Inspector 2) y un checklist estructurado por dimensión de calidad
 ambigüedad, completitud, consistencia, verificabilidad y factibilidad, lo que produce
una evidencia auditable de que la validación ocurrió y de qué se encontró, en lugar de una
declaración de que ``se revisó'' el documento sin registro verificable.

En el subgrupo de origen de dos de los cinco integrantes actuales del equipo, la inspección
se ejecutó sobre el ERS del PFC con dos inspectoras que trabajaron de forma independiente
antes de la reunión de consolidación: una desde la perspectiva de consistencia y otra desde
las perspectivas de ambigüedad, verificabilidad, completitud y factibilidad. El resultado
fue de 33 defectos
reportados en total, de los cuales 27 se acreditaron como únicos tras descartar 6
coincidencias de frontera detectadas por ambas inspectoras. De esos 27 defectos, 22 se
corrigieron directamente sobre el ERS y quedaron reflejados en la versión 1.1 del
documento, mientras que 5 defectos menores permanecieron abiertos y fueron reconocidos como
tales en el informe, en lugar de declararse cerrados sin respaldo. 

Los hallazgos que excedían una corrección directa por afectar el alcance, la calidad de un requisito no
funcional o una restricción normativa se tramitaron mediante \textbf{solicitudes de
	cambio (RFC)} y se sometieron a un \textbf{Comité de Control de Cambios (CCB)} simulado,
conforme al proceso de gestión de cambios que describe ISO/IEC/IEEE 12207:2017
\cite{iso12207} y que formaliza el PMBOK 7.\textsuperscript{a} edición \cite{pmbok7}: cada
RFC se registró, se analizó su impacto y se sometió a una decisión motivada del comité
aprobada, aprobada con condiciones o rechazada, dejando acta. Los otros tres subgrupos
del equipo ejecutaron el mismo procedimiento sobre sus propias copias del ERS, con
resultados cuantitativos distintos pero metodológicamente equivalentes: entre 18 y 23
defectos únicos según el subgrupo, siempre con roles Fagan diferenciados y siempre con al
menos tres RFC tramitadas ante un CCB.

El artefacto que produjo la PE4 fue, en cada subgrupo, una versión corregida y con línea
base del ERS, acompañada del registro de defectos, del acta de inspección y de las
solicitudes de cambio resueltas. Este es el punto en el que el proceso de las cuatro
prácticas individuales se detiene: lo que la PE4 entregó no fue todavía el ERS único del
equipo de la PE5, sino cuatro ERS paralelos, cada uno validado de forma independiente sobre
el mismo dominio de problema.

\subsection{PE5: integración, cierre de cambios pendientes y defensa}

La PE5 no repite ninguna de las cuatro prácticas anteriores: las integra. El primer paso
metodológico del equipo, ya reunido bajo un mismo repositorio, fue reconciliar las cuatro
versiones del ERS surgidas de PE1--PE4 en un documento único y consolidar en él los
requisitos, el modelo de datos, el modelo funcional, el modelo de comportamiento y la
matriz de trazabilidad de los cuatro subgrupos. Como parte de ese cierre, el equipo revisó
las solicitudes de cambio que habían quedado pendientes de las cuatro inspecciones de la
PE4 y las resolvió formalmente ante un nuevo CCB: nueve RFC fueron decididas sin dejar
ninguna abierta, lo que motivó la corrección de varios requisitos funcionales y no
funcionales del ERS consolidado y, en los casos en que la evidencia disponible no permitía
aceptar el cambio propuesto, el rechazo justificado de la solicitud con su registro
correspondiente como limitación de alcance reconocida. 

Sobre esa base consolidada y ya sin
cambios pendientes, la PE5 añadió lo que ninguna de las cuatro prácticas anteriores había
cubierto: la auditoría de calidad del ERS con las seis métricas de
la sección de métricas de calidad, el cierre definitivo de la trazabilidad end-to-end de la
sección de trazabilidad y matriz final, la especificación de los componentes de
inteligencia artificial de la sección de requisitos de IA, y la preparación de la defensa técnica individual
ante el tribunal. En ese sentido, la PE5 es menos una sexta técnica que una disciplina de
cierre: aplica a la totalidad del proceso anterior los mismos criterios de verificabilidad
que EARS exigió a cada requisito individual en la PE3, pero ahora a la escala del proyecto
integrador completo.

\newpage

\subsection{Tabla resumen del proceso}

La Tabla~\ref{tab:resumen-metodologia} sintetiza la actividad, la técnica, el artefacto
producido y la norma de referencia de cada práctica, en el mismo orden en que se
describieron en esta sección.

\begin{table}[H]
	\centering
	\caption{Resumen del proceso de Ingeniería de Requisitos aplicado a MundiPets, PE1--PE5.}
	\label{tab:resumen-metodologia}
	\footnotesize
	\renewcommand{\arraystretch}{1.3}
	\begin{tabularx}{\textwidth}{
			|>{\raggedright\arraybackslash}p{0.9cm}
			|>{\raggedright\arraybackslash}p{2.3cm}
			|>{\raggedright\arraybackslash}X
			|>{\raggedright\arraybackslash}p{3.1cm}
			|>{\raggedright\arraybackslash}p{3.4cm}|}
		\hline
		\textbf{Práctica} & \textbf{Actividad} & \textbf{Técnica} & \textbf{Artefacto} & \textbf{Norma} \\
		\hline
		PE1 & Caracterización del sistema y su contexto & Matriz poder/interés; ocho perspectivas de contexto; declaración de visión (Moore) & Mapa de restricciones y necesidades preliminares & PMBOK 7.\textsuperscript{a} ed.; Pohl \\
		\hline
		PE2 & Elicitación, análisis y priorización & Entrevista estructurada; cuestionario digital; brainstorming; MoSCoW; Kano & Lista de requisitos depurada, clasificada y priorizada & ISO/IEC/IEEE 29148:2018; IREB CPRE FL; ISO/IEC 25010:2023 \\
		\hline
		PE3 & Especificación y modelado & Sintaxis EARS; modelo de datos (ER/clases UML); modelo funcional (DFD/actividades); modelo de comportamiento (estados/secuencia) & ERS corregido, modelado y con matriz de trazabilidad preliminar & ISO/IEC/IEEE 29148:2018; SWEBOK v4.0 \\
		\hline
		PE4 & Validación y gestión de cambios & Inspección formal Fagan; RFC; CCB simulado & ERS corregido, con línea base y registro de defectos & ISO/IEC/IEEE 29148:2018; ISO/IEC/IEEE 12207:2017; PMBOK 7.\textsuperscript{a} ed. \\
		\hline
		PE5 & Integración, auditoría y defensa & Consolidación de las cuatro versiones; cierre de RFC pendientes; métricas ISO/IEC 25010; trazabilidad end-to-end; requisitos de IA & ERS final único, matriz de trazabilidad final, informe integrador & ISO/IEC/IEEE 29148:2018; ISO/IEC 25010:2023; Reglamento (UE) 2024/1689 \\
		\hline
	\end{tabularx}
\end{table}

\subsection{Fundamento normativo del proceso seguido}

El proceso descrito no se apoya en una sola fuente sino en un conjunto de normas y marcos
que se complementan sin superponerse. La ISO/IEC/IEEE 29148:2018 encuadra el proceso de punta a punta (elicitación, análisis, especificación, validación y gestión) y es la
norma contra la cual se verifica que cada artefacto producido en PE1--PE5 sea,
formalmente, parte de un proceso reconocido de Ingeniería de Requisitos y no una secuencia
de tareas de curso sin relación entre sí \cite{ieee29148}. La ISO/IEC/IEEE 15288:2023 y la
ISO/IEC/IEEE 12207:2017 sitúan ese proceso dentro del ciclo de vida del sistema y del
software respectivamente, y son las que sustentan formalmente la actividad de gestión de
cambios ejecutada en la PE4 y cerrada en la PE5 \cite{iso15288,iso12207}. Pohl aporta el
fundamento conceptual de las actividades que preceden a la especificación formal (elicitación, análisis, negociación) y es la fuente directa de las técnicas aplicadas en
la PE1 y en la priorización de la PE2 \cite{pohl2025}. El SWEBOK v4.0 enmarca estas
actividades dentro del cuerpo de conocimiento reconocido de la ingeniería de software y
respalda en particular la práctica de modelado de la PE3 \cite{swebok2024}. El IREB CPRE FL
v3.3.0 aporta el catálogo verificado de técnicas de elicitación (entrevista, cuestionario,
brainstorming) que sustenta la selección hecha en la PE2 \cite{ireb2026}. Ninguna de estas
fuentes sustituye a las demás: cada una responde a una etapa distinta del mismo proceso, y
es esa complementariedad, y no la acumulación de citas, la que permite sostener que las
decisiones metodológicas del equipo en PE1--PE5 no fueron arbitrarias.

\section{Referencias}

\printbibliography[heading=none]

\end{document}
